% Ad Atticum 15.15 --- headnote
% ad-atticum-15-15, 13 June 44 BC, in Antiati

\textit{Cicero to Atticus, written at Antium on 13 June 44 BC ---
Perseus dateline \textit{Scr. in Antiat\textsubscript{i} Id. Iun.
a. 710 (44)}. Cicero is on the move, two months after the
assassination of Caesar and at the point when Brutus and Cassius
are reconsidering whether to remain in Italy at all. The
Buthrotian affair (see 15.14) is still in play: Lucius Antonius,
the consul's brother, is reported to be agitating against the
city, and Cicero has prepared a sworn deposition to be sealed at
Atticus's word. There follow practical instructions about a sum
owed to the Arpinate community, which is to be paid to the
aedile Lucius Fadius and no one else.}

\textit{Section 2 contains the famous outburst, \textit{reginam
odi}: ``I detest the queen.'' Cleopatra, who had been resident
in Caesar's gardens across the Tiber from her arrival in 46 down
to her flight after the Ides, had apparently promised Cicero
something --- something literary, evidently, and worth being
proud of in public --- and had failed to deliver. Her agents
Ammonius and Sara have compounded the offence with personal
insolence. The bitter little aside ``they suppose I have no
spirit, or rather no stomach'' is the kind of pun (\textit{animus}
``courage'' versus \textit{stomachus} ``spleen'') that survives
imperfectly in any language. The letter closes with two
financial nuisances: Eros, Cicero's freedman accountant, has
left Cicero short of cash to fund his projected travel, and
young Marcus Cicero, studying in Athens, has gone without his
allowance since April --- an embarrassment that pains Cicero
the more because the boy is too modest to complain of it to his
father directly, and has told Tiro instead.}
